\usepackage{circuitikz}
\usepackage{siunitx}
\usepackage{pgffor}
\usepackage{tikz}
\usetikzlibrary{intersections}
\usepackage{tkz-euclide}
\usetikzlibrary{arrows.meta}
\usepackage{ifthen}
\def\LW{1pt}
\tikzset{every path/.style={line width=\LW}}
\tikzset{arrows/.style={line width=1pt}}

\newcommand\relay[1] {%1 start coords
    \draw #1 + (1.5,2) -- ++(1.5, 1.25) -- ++(-.375,-.25) to[L] ++(0,-1) -- ++(-1.125,0);
    % \draw #1 -- ++(1,0) to[L] ++(0,1) -- ++(.5,.5) -- ++(0,.5);
    \draw #1 + (.5, -.5) rectangle ++(2.5, 1.5);
    \draw #1 + (3,0) -- ++(1.875,0) to[nos] ++(0,1) -- ++(1.125,0);
}

% Orbit is a Vici rotary valve with n ports and 2 positions
\newcommand{\orbit}[5]{% #1 = name, #2 = start coords, #3 = n, #4 = scale, #5 = 0/1 for off/on
\draw #2 circle (1.5*#4);
\coordinate (#1) at #2;
\draw #2 node[]{#1};
\def\step{\the\numexpr(360/#3)}
\foreach \n in {1,...,{\the\numexpr#3/2}}{
    \def\ct{\the\numexpr(2*\n-2)}
    \coordinate (a) at ({#4*sin(\ct*\step)}, {#4*cos(\ct*\step});
    \ifodd#5
    \coordinate (b) at ({#4*sin((\ct+1)*\step)},{#4*cos((\ct+1)*\step)});
    \else
    \coordinate (b) at ({#4*sin((\ct-1)*\step)},{#4*cos((\ct-1)*\step)});
    \fi
    \draw ($#2 + 2*(a)$) 
    -- ($#2 + (a)$) 
    -- ($#2 + (b)$) 
    -- ($#2 + 2*(b)$);
    \draw ($#2 + (a)$) circle (#4*.125);
    \draw ($#2 + (b)$) circle (#4*.125);
}
\foreach \n in {1, ...,#3}{
    \coordinate (a) at ({#4*sin((\n-1)*\step)}, {#4*cos((\n-1)*\step)});
    \coordinate (#1-\n) at ($#2 + 2*(a)$);
    \draw ($#2 + .675*(a)$) node[]{\n};
}
}

\newcommand{\cross}[3]{%1 start coords, %2 = stop coords, %3 crossing line
  % find intersection of first and second line
  \path [name path=line 2] (#2) -- (#1);
  \path [name intersections={of = #3 and line 2}];
  \coordinate (S)  at (intersection-1);
  % path a circle around this intersection for the arc
  \path[name path=circle] (S) circle(3pt);
  % find intersections of second line and circle
  \path [name intersections={of = circle and line 2}];
  \coordinate (IA)  at (intersection-1);
  \coordinate (IB)  at (intersection-2);
  % Draw the crossing
  \draw (#1) -- (IA) to[crossing, bipoles/crossing/size=.4] (IB) -- (#2);
}

% TWV = Three-Way Valve
\newcommand{\TWV}[5]{% #1 = name, #2 =start coords, #3 = size, #4 = rotation, #5 = 1/0 for NC/NO
    \coordinate (#1-C) at #2;
    \coordinate (center) at ($#2 + ({1.25*#3*sin(#4)},{1.25*#3*cos(#4)})$);
    \def\y{5*#3*cos(#4)}
    \def\x{5*#3*sin(#4)}
    \draw (center) circle (#3) node[yshift=\y, xshift=\x]{#1};
    
    \coordinate (#1-NC) at ($(center) + ({1.25*#3*cos(#4)},{-1.25*#3*sin(#4)})$);
    \coordinate (#1-NO) at ($(center) + ({-1.25*#3*cos(#4)},{1.25*#3*sin(#4)})$);
    \ifodd#5
    \draw (#1-C) -- ($#2 + ({.25*#3*sin(#4)}, {.25*#3*cos(#4)})$) arc ((#4-180):(#4-270):-#3 and #3) -- (#1-NO);
    \draw (#1-NC) -- ($(center) + ({#3*cos(#4)},{-1*#3*sin(#4)})$);
    \else
    \draw (#1-C) -- ($#2 + ({.25*#3*sin(#4)}, {.25*#3*cos(#4)})$) arc ((180-#4):(90-#4):#3 and #3) -- (#1-NC);
    \draw (#1-NO) -- ($(center) + ({-1*#3*cos(#4)},{#3*sin(#4)})$);
    \fi
}

\newcommand{\mfc}[3]{% #1 = name, #2 = start coords, #3 = size
    \coordinate (#1-inlet) at #2;
    \coordinate (#1-outlet) at ($#2 + (#3*2.5,0)$);
    \draw #2 -- ++(0.25*#3,0);
    \draw #2 + (.25*#3,#3*.5) rectangle ++(#3*2.25,-#3*.5) node[pos=.5]{#1};
    \ifnum 0>#3
    \draw #2 + ({.75*#3}, {-1*#3}) rectangle ++(#3*1.75, {-1*#3*.5});
    \else
    \draw #2 + ({.75*#3}, {#3}) rectangle ++(#3*1.75, {#3*.5});
    \fi
    \draw #2 + (2.25*#3,0) -- (#1-outlet);
}

\newcommand{\block}[3]{% #1 = name, #2 = start coords, #3 = size
    \coordinate (#1-inlet) at #2;
    \coordinate (#1-outlet) at ($#2 + (#3*2.5,0)$);
    \draw #2 -- ++(0.25*#3,0);
    \draw #2 + (.25*#3,#3*.75) rectangle ++(#3*2.25,-#3*.75) node[pos=.5]{#1};
    \draw #2 + (2.25*#3,0) -- (#1-outlet);
    % \draw #2 + (1.5*#3,0) node[right, text width=\txtwidth]{#1};
}

\newcommand{\pump}[4]{% #1 = name, #2 = start coords, #3 = size, #4 = rotation
    \coordinate (#1-inlet) at #2;
    \coordinate (center) at ($#2 + ({1.25*#3*sin(#4)},{1.25*#3*cos(#4)})$);
    \coordinate (#1-outlet) at ($(center) + ({1.25*#3*sin(#4)},{1.25*#3*cos(#4)})$);
    \draw (center) circle (#3) node[]{#1};
    \draw (#1-inlet) -- ($#2 + ({.25*#3*sin(#4)}, {.25*#3*cos(#4)})$);
    \draw ($(center) + ({#3*sin(#4)}, {#3*cos(#4)})$) -- (#1-outlet);
    % Pump angled lines
    \draw ($(center) + ({#3*cos(50-#4)}, {#3*sin(50-#4)})$) --
    ($(center) + ({#3*cos(-75-#4)}, {#3*sin(-75-#4)})$);
    \draw ($(center) + ({#3*cos(50-#4)}, {-#3*sin(50-#4)})$) --
    ($(center) + ({#3*cos(-75-#4)}, {-#3*sin(-75-#4)})$);

}

\newcommand{\gc}[4]{% #1 = name, #2 = start coords, #3 = size, #4 = detector
    \draw ($#2 + (-.375*#3,0)$) rectangle ++(.75*#3, -.5*#3);
    \draw ($#2 + (-.375*#3,-.5*#3)$) rectangle ++(2.5*#3, -2.5*#3);
    \draw ($#2 + (1.375*#3,-.5*#3)$) rectangle ++(.75*#3, .5*#3) node[pos=.5]{#4};
    \draw ($#2 + (0.875*#3,-3*#3)$) node[above]{#1};
    % Column Coils
    \draw #2 -- ++(0, -1.5*#3);
    \draw ($#2 + (.75*#3,-1.5*#3)$) circle (.75*#3);
    \draw ($#2 + (.875*#3,-1.5*#3)$) circle (.75*#3);
    \draw ($#2 + (1*#3,-1.5*#3)$) circle (.75*#3);
    \draw ($#2 + (1.75*#3,-.5*#3)$) -- ++(0, -1*#3);
}


\newcommand{\mpv}[5]{% #1 = name, #2 = start coords, #3 = size, #4 = rotation, #5 = ports
    \coordinate (#1-outlet) at #2;
    \coordinate (center) at ($#2 + 1.25*#3*({cos(#4)},{sin(#4)})$);
    \draw #2 -- ++($.2*($(center) - #2$)$);
    \draw (center) circle (#3) node[]{#1};
    % \ifthenelse{#5>7}{\def\N{7}}{\def\N{#5}}
    
    \foreach \x in {1, ..., #5}{
            \def\dT{120/#5}
            \def\T{(\x-0.5)*\dT - 60}
            This number is not zero.
            \coordinate (start) at ($(center) + #3*({cos(\T+#4)},{sin(\T+#4)})$);
            \coordinate (extend) at ($0.75*#3*($(start) - (center)$)$);
            \def\mag{#3*(2-1.75*cos(\T))}
            \coordinate (reach) at (${\mag}*({cos(#4)}, {sin(#4)})$);
            \draw (start) -- ++(extend) -- ++(reach);
            \coordinate (#1-\x) at ($(start) + (extend) + (reach)$);
        }
}

\newcommand\valve[4]{% #1 = name, #2 = start coords, #3 = size, #4 = rotation
    \coordinate (#1-in) at #2;
    \coordinate (A) at ($#3*({cos(#4)}, {sin(#4)})$);
    \coordinate (B) at ($#3*({cos(#4+90)}, {sin(#4+90)})$);
    \coordinate (#1-out) at ($(#1-in) + (A)$);
    \draw (#1-in) -- ++($0.25*#3*(A)$) 
    -- ++($0.125*(B)$) 
    --($(#1-in) + .75*(A) - .125*(B)$)
    -- ++($0.25*(B)$)
    -- ($(#1-in)+.25*(A) - .125*(B)$)
    -- ++($0.125*(B)$);
    \draw (#1-out) -- ++($-.25*#3*(A)$);
}

\newcommand\rotameter[4]{% #1 = name, #2 = start coords, #3 = size, #4 = rotation
    \coordinate (#1-in) at #2;
    \def\r{#3*.75*.5}
    \def\x{#3*.25}
    \pgfmathsetmacro{\y}{\r*sin(90-asin(\r/(\r+\x)))}
    \pgfmathsetmacro{\xp}{(\x+\r)-\r*cos(90-asin(\r/(\r+\x)))}
    \coordinate (A) at ($#3*({cos(#4)}, {sin(#4)})$);
    \coordinate (B) at ($#3*({cos(#4+90)}, {sin(#4+90)})$);
    \coordinate (#1-out) at ($(#1-in) + {\x+2*\r}*(A)$);
    \coordinate (start) at ($#2 + \xp*(A) + \y*(B)$);
    \coordinate (stop) at ($#2 + \xp*(A) - \y*(B)$);
    \coordinate (#1-center) at ($#2 + {\x+\r}*(A)$);
    \draw (start) -- (#1-in) -- (stop);
    \draw (#1-center) circle (#3*\r);
}
